\section{Class P --- Binary Arith --- Add Sub}
%%% generated by the blueprint pipeline from TCSlib.Complexity.ClassP.BinaryArith.AddSub
\subsection*{Class P --- Binary Arith --- Add Sub: Introduction}

Declarations of the Lean module \texttt{TCSlib.Complexity.ClassP.BinaryArith.AddSub}.

\subsection*{Class P --- Binary Arith --- Add Sub: Results}

\begin{definition}[Boolean Carry Function]
\lean{Complexity.BinArith.carry}
\label{Complexity.BinArith.carry}
\leanok
For Boolean bits $c$, $x$, and $y$, define their carry output to be true exactly when
either both $x$ and $y$ are true, or $c$ is true and exactly one of $x$ and $y$ is true.
\end{definition}

\begin{definition}[Binary sum bit]
\lean{Complexity.BinArith.sumBit}
\label{Complexity.BinArith.sumBit}
\leanok
For Boolean bits $c$, $x$, and $y$, define their sum bit as the exclusive-or of all
three: $(x\mathbin{\mathrm{xor}} y)\mathbin{\mathrm{xor}} c$.
\end{definition}

\begin{definition}[Binary Borrow Bit]
\lean{Complexity.BinArith.borrow}
\label{Complexity.BinArith.borrow}
\leanok
For input bits $c,x,y\in\{0,1\}$, define the borrow output to be $1$ exactly when
$(1-x)y+(1-(x\mathbin{\mathrm{xor}}y))c$ is nonzero; otherwise it is $0$.
\end{definition}

\begin{theorem}[Serial Binary Addition]
\lean{Complexity.BinArith.serial_add}
\label{Complexity.BinArith.serial_add}
\leanok
\uses{Complexity.BinArith.bitsVal_eq_headD, Complexity.BinArith.carry, Complexity.BinArith.padZip, Complexity.BinArith.serialOut, Complexity.BinArith.serialState, Complexity.BinArith.sumBit, Complexity.LogProg.bitsVal}
\difficulty{4}
For every natural number $N$, bit strings $a,b$ with $|a|,|b|\leq N$, and initial carry
bit $c$, process the first $N$ aligned positions of $a$ and $b$, padding either word by
zeros after it ends. At each position output the sum bit and update the carry by the
usual binary addition rules. The value of the resulting output word plus $2^N$ times the
final carry equals $\langle a\rangle+\langle b\rangle+c$, where $c$ on the right is
interpreted as $0$ or $1$.
\end{theorem}

\begin{theorem}[Value identity for serial subtraction]
\lean{Complexity.BinArith.serial_sub}
\label{Complexity.BinArith.serial_sub}
\leanok
\uses{Complexity.BinArith.bitsVal_eq_headD, Complexity.BinArith.borrow, Complexity.BinArith.padZip, Complexity.BinArith.serialOut, Complexity.BinArith.serialState, Complexity.BinArith.sumBit, Complexity.LogProg.bitsVal}
\difficulty{4}
Let $N$ be a natural number, let $a,b$ be finite bit strings of lengths at most $N$, and
let $c$ be an initial borrow bit. Pad $a$ and $b$ with zero bits to $N$ positions, then
process the positions from least to most significant using the subtraction output bit
and borrow update. If $o$ is the resulting $N$-bit output and $c_N$ the final borrow
bit, then $\operatorname{val}(o)+\operatorname{val}(b)+c=\operatorname{val}(a)+2^N c_N$,
where bits are interpreted as $0$ or $1$.
\end{theorem}

\begin{theorem}[Bounds for decoded binary values]
\lean{Complexity.BinArith.serial_bounds}
\label{Complexity.BinArith.serial_bounds}
\leanok
\uses{Complexity.LogProg.bitsVal, Complexity.LogProg.bitsVal_lt_two_pow, Complexity.length_pairFstD_le, Complexity.length_pairSndD_le, Complexity.pairFstD, Complexity.pairSndD}
\difficulty{3}
For every finite bit string $z$, let $a$ and $b$ be its decoded first and second
components, with either component taken to be the empty string when decoding fails.
Interpreting a bit string as a binary numeral with its first bit least significant, both
values satisfy $\operatorname{val}(a)<2^{|z|+2}$ and $\operatorname{val}(b)<2^{|z|+2}$.
\end{theorem}

\begin{theorem}[Correctness of serial subtraction]
\lean{Complexity.BinArith.serial_sub_spec}
\label{Complexity.BinArith.serial_sub_spec}
\leanok
\uses{Complexity.BinArith.borrow, Complexity.BinArith.length_padZip, Complexity.BinArith.length_serialOut, Complexity.BinArith.padZip, Complexity.BinArith.serialOut, Complexity.BinArith.serialState, Complexity.BinArith.serial_bounds, Complexity.BinArith.serial_sub, Complexity.BinArith.sumBit, Complexity.LogProg.bitsVal, Complexity.LogProg.bitsVal_lt_two_pow, Complexity.length_pairFstD_le, Complexity.length_pairSndD_le, Complexity.pairFstD, Complexity.pairSndD}
\difficulty{5}
For any finite bit string $z$, let $a$ and $b$ be its decoded components, taking a
failed decoding to give empty components. Run the $|z|+2$-position serial subtraction
procedure on these strings, padded with zeros and initialized with borrow $0$. Its final
borrow bit is $1$ exactly when $\operatorname{val}(a)<\operatorname{val}(b)$. If
$\operatorname{val}(b)\le\operatorname{val}(a)$, the output has value
$\operatorname{val}(a)-\operatorname{val}(b)$.
\end{theorem}

\begin{definition}[Binary addition of encoded inputs]
\lean{Complexity.BinArith.add}
\label{Complexity.BinArith.add}
\leanok
\uses{Complexity.LogProg.bitsVal, Complexity.pairFstD, Complexity.pairSndD}
For a Boolean list $z$, define its sum to be the binary expansion of the sum of the
natural numbers represented by the two components obtained by splitting and decoding
$z$.
\end{definition}

\begin{definition}[Binary subtraction of encoded inputs]
\lean{Complexity.BinArith.sub}
\label{Complexity.BinArith.sub}
\leanok
\uses{Complexity.LogProg.bitsVal, Complexity.pairFstD, Complexity.pairSndD}
For a Boolean list $z$, define its difference to be the binary expansion of the
truncated natural-number subtraction of the value represented by the second decoded
component of $z$ from the value represented by the first decoded component.
\end{definition}

\begin{theorem}[Addition on encoded natural numbers]
\lean{Complexity.BinArith.add_pairEncode}
\label{Complexity.BinArith.add_pairEncode}
\leanok
\uses{Complexity.BinArith.add, Complexity.LogProg.bitsVal_bits, Turing.pairEncode}
\difficulty{1}
For any natural numbers $a$ and $b$, decoding their binary representations from the
self-delimiting pair encoding and applying the binary addition operation returns the
binary representation of $a+b$.
\end{theorem}

\begin{theorem}[Subtraction on encoded binary inputs]
\lean{Complexity.BinArith.sub_pairEncode}
\label{Complexity.BinArith.sub_pairEncode}
\leanok
\uses{Complexity.BinArith.sub, Complexity.LogProg.bitsVal_bits, Turing.pairEncode}
\difficulty{1}
For any natural numbers $a$ and $b$, encode their binary expansions as a pair by
doubling every bit of the first expansion, appending the separator $01$, and then the
second expansion. Apply the decoded subtraction operation to this encoded string. The
result is the binary expansion of $a-b$, where subtraction on natural numbers is
truncated at zero.
\end{theorem}

\begin{theorem}[Polynomial-time binary addition]
\lean{Complexity.BinArith.polyTimeComputable_add}
\label{Complexity.BinArith.polyTimeComputable_add}
\leanok
\uses{Complexity.BinArith.add, Complexity.BinArith.carry, Complexity.BinArith.padZip, Complexity.BinArith.polyTimeComputable_normBits, Complexity.BinArith.polyTimeComputable_serial, Complexity.BinArith.serialState, Complexity.BinArith.serial_add, Complexity.BinArith.sumBit, Complexity.LogProg.bitsVal, Complexity.LogProg.bitsVal_lt_two_pow, Complexity.PolyTimeComputable, Complexity.PolyTimeComputable.comp, Complexity.length_pairFstD_le, Complexity.length_pairSndD_le, Complexity.pairFstD, Complexity.pairSndD}
\difficulty{6}
The function on bit strings that decodes the two components of a pair, interprets each
component as a natural number with the first bit least significant, adds those numbers,
and returns the binary expansion of their sum is polynomial-time computable.
\end{theorem}

\begin{definition}[Binary less-than comparison]
\lean{Complexity.BinArith.lt}
\label{Complexity.BinArith.lt}
\leanok
\uses{Complexity.LogProg.bitsVal, Complexity.pairFstD, Complexity.pairSndD}
For a list of Boolean values encoding a pair of binary numbers, return `true` exactly
when the number encoded by the first component is strictly less than the number encoded
by the second.
\end{definition}

\begin{definition}[Binary less-than-or-equal comparison]
\lean{Complexity.BinArith.le}
\label{Complexity.BinArith.le}
\leanok
\uses{Complexity.LogProg.bitsVal, Complexity.pairFstD, Complexity.pairSndD}
For a list of Boolean values encoding a pair of binary numbers, return `true` exactly
when the number encoded by the first component is less than or equal to the number
encoded by the second.
\end{definition}

\begin{definition}[Binary equality comparison]
\lean{Complexity.BinArith.beq}
\label{Complexity.BinArith.beq}
\leanok
\uses{Complexity.LogProg.bitsVal, Complexity.pairFstD, Complexity.pairSndD}
For a list of Boolean values encoding a pair of binary numbers, return `true` exactly
when the numbers encoded by the two components are equal.
\end{definition}

\begin{theorem}[Comparison of encoded integers]
\lean{Complexity.BinArith.lt_pairEncode}
\label{Complexity.BinArith.lt_pairEncode}
\leanok
\uses{Complexity.BinArith.lt, Complexity.LogProg.bitsVal_bits, Turing.pairEncode}
\difficulty{1}
For natural numbers $a$ and $b$, encode their binary expansions as a pair. The
comparison of the values decoded from the two components of this pair is true exactly
when $a<b$.
\end{theorem}

\begin{theorem}[Comparison on encoded binary numerals]
\lean{Complexity.BinArith.le_pairEncode}
\label{Complexity.BinArith.le_pairEncode}
\leanok
\uses{Complexity.BinArith.le, Complexity.LogProg.bitsVal_bits, Turing.pairEncode}
\difficulty{1}
For all natural numbers $a$ and $b$, the comparison predicate applied to the
self-delimiting encoding of their binary expansions is true exactly when $a \le b$.
\end{theorem}

\begin{theorem}[Equality test on encoded natural numbers]
\lean{Complexity.BinArith.beq_pairEncode}
\label{Complexity.BinArith.beq_pairEncode}
\leanok
\uses{Complexity.BinArith.beq, Complexity.LogProg.bitsVal_bits, Turing.pairEncode}
\difficulty{1}
For any natural numbers $a$ and $b$, the Boolean equality test applied to the
self-delimiting encoding of their binary representations returns true exactly when
$a=b$, and false otherwise.
\end{theorem}

\begin{theorem}[Polynomial-Time Comparison of Encoded Integers]
\lean{Complexity.BinArith.polyTimeComputable_lt}
\label{Complexity.BinArith.polyTimeComputable_lt}
\leanok
\uses{Complexity.BinArith.borrow, Complexity.BinArith.getLastD_serialOut_self, Complexity.BinArith.length_padZip, Complexity.BinArith.lt, Complexity.BinArith.polyTimeComputable_lastBit, Complexity.BinArith.polyTimeComputable_serial, Complexity.BinArith.serial_sub_spec, Complexity.PolyTimeComputable, Complexity.PolyTimeComputable.comp}
\difficulty{3}
Encode a pair of bit strings by doubling each bit of the first string, appending the
separator $01$, and then appending the second string. Parse an input at aligned pairs
from the left, treating $00$ and $11$ as doubled bits of the first string and the first
$01$ as the separator; if parsing fails, take both decoded strings to be empty.
Interpret each decoded string as a binary numeral with its first bit least significant.
The function that returns the one-bit string $1$ when the first decoded value is
strictly less than the second, and $0$ otherwise, is computable in polynomial time.
\end{theorem}

\begin{theorem}[Polynomial-Time Comparison]
\lean{Complexity.BinArith.polyTimeComputable_le}
\label{Complexity.BinArith.polyTimeComputable_le}
\leanok
\uses{Complexity.BinArith.le, Complexity.BinArith.lt, Complexity.BinArith.polyTimeComputable_lt, Complexity.LogProg.bitsVal, Complexity.PolyTimeComputable, Complexity.PolyTimeComputable.comp, Complexity.pairFstD, Complexity.pairSndD, Complexity.polyTimeComputable_not, Complexity.polyTimeComputable_pairSwap}
\difficulty{4}
The function that sends a finite Boolean word $z$ to the one-bit word $[\mathsf{true}]$
if the value of its first decoded component is at most the value of its second decoded
component, and to $[\mathsf{false}]$ otherwise, is polynomial-time computable.
\end{theorem}

\begin{theorem}[Polynomial-time equality test]
\lean{Complexity.BinArith.polyTimeComputable_beq}
\label{Complexity.BinArith.polyTimeComputable_beq}
\leanok
\uses{Complexity.BinArith.beq, Complexity.BinArith.le, Complexity.BinArith.polyTimeComputable_le, Complexity.LogProg.bitsVal, Complexity.PolyTimeComputable, Complexity.PolyTimeComputable.comp, Complexity.pairFstD, Complexity.pairFstD_pairEncode, Complexity.pairSndD, Complexity.pairSndD_pairEncode, Complexity.polyTimeComputable_and, Complexity.polyTimeComputable_pairSwap}
\difficulty{5}
The function on bit strings that decodes the two components of a pair, interprets each
as a natural number with the first bit least significant, and returns the one-bit word
$[1]$ when the values are equal and $[0]$ otherwise, is polynomial-time computable.
\end{theorem}

\begin{theorem}[Polynomial-time binary subtraction]
\lean{Complexity.BinArith.polyTimeComputable_sub}
\label{Complexity.BinArith.polyTimeComputable_sub}
\leanok
\uses{Complexity.BinArith.borrow, Complexity.BinArith.lt, Complexity.BinArith.polyTimeComputable_lt, Complexity.BinArith.polyTimeComputable_normBits, Complexity.BinArith.polyTimeComputable_serial, Complexity.BinArith.serial_sub_spec, Complexity.BinArith.sub, Complexity.BinArith.sumBit, Complexity.LogProg.bitsVal, Complexity.PolyTimeComputable, Complexity.PolyTimeComputable.comp, Complexity.pairFstD, Complexity.pairSndD, Complexity.polyTimeComputable_const, Complexity.polyTimeComputable_ite}
\difficulty{4}
The function that takes a bit string, applies the total first and second projections,
interprets each resulting bit string as a natural number with least significant bit
first, subtracts the second value from the first using truncated subtraction, and
returns the binary expansion of the result is polynomial-time computable.
\end{theorem}

\begin{theorem}[Polynomial-time decidability of comparison]
\lean{Complexity.BinArith.lt_mem_P}
\label{Complexity.BinArith.lt_mem_P}
\leanok
\uses{Complexity.BinArith.lt, Complexity.BinArith.polyTimeComputable_lt, Complexity.P, Complexity.mem_P_of_test}
\difficulty{1}
Let $z$ be a finite binary string, and let $a$ and $b$ be the two binary strings
obtained by parsing it as a pair encoding, using the empty string for either projection
when the parse is invalid. Consider the language of strings $z$ for which the natural
number represented by $a$ is strictly less than the natural number represented by $b$;
this language is decidable in polynomial time.
\end{theorem}

\begin{theorem}[Polynomial-time binary comparison language]
\lean{Complexity.BinArith.le_mem_P}
\label{Complexity.BinArith.le_mem_P}
\leanok
\uses{Complexity.BinArith.le, Complexity.BinArith.polyTimeComputable_le, Complexity.P, Complexity.mem_P_of_test}
\difficulty{1}
For a bit string $z$, let $a$ and $b$ be the two components parsed from $z$, taking both
to be the empty string if parsing fails, and interpret each bit string as a binary
numeral with its first bit least significant. The language of strings for which the
value of $a$ is at most the value of $b$ belongs to $P$.
\end{theorem}

\begin{theorem}[Polynomial-time decidability of encoded equality]
\lean{Complexity.BinArith.beq_mem_P}
\label{Complexity.BinArith.beq_mem_P}
\leanok
\uses{Complexity.BinArith.beq, Complexity.BinArith.polyTimeComputable_beq, Complexity.P, Complexity.mem_P_of_test}
\difficulty{1}
The language of bit strings whose two decoded components have equal binary values is in
$\mathrm{P}$. Here decoding a bit string that is not a valid pair encoding gives the
empty word for each component.
\end{theorem}
